\section{Application: Criticality and Trade-Disruption Analysis}
\label{sec:application}

The MEKH is built to be used. This section gives the two analyses the institute's analysts run on it: which materials are structurally critical to a mineral's downstream reach, and what a concentrated exporter's withdrawal would remove.

\subsection{Criticality by reachability}
Each \tool{} is constructed for a base material, e.g., boron, which we denote $m_{base}$; the set of materials that $m_{base}$, in combination with other raw materials, can be transformed into is its \textit{reach}. The criticality of any other material is the reduction in the reach of $m_{base}$ when that material is removed; from a policy perspective, removal could correspond to banning the production and use of that material, and criticality quantifies the cascading effect.

\begin{definition}
\label{def:pre}
\textbf{Precursor Set.}
Given a \tool{} $\Hcal(M, E)$, where each hyperedge $e \in E$ has a source set $e_s$ and a target set $e_t$, we define the precursor set $\mathtt{Pre} = \bigcup_e e_s - \bigcup_e e_t$ to be the set of materials which appear only in hyperedge sources.
\end{definition}

\begin{definition}
\label{def:reach}
\textbf{Reachable set.}
Given a \tool{} $\Hcal(M, E)$ and a set $m_{start} \subseteq M$, a material $m_i$ is considered reachable if it is in the
target set of a hyperedge whose source set contains only materials that are in $m_{start}$ or are themselves reachable. The 
reachable set, $R(m_{start}, \Hcal)$ is the set of all nodes that are reachable from $m_{start}$ in $\Hcal$.
\end{definition}

\begin{definition}
\label{def:criticality}
\textbf{Criticality.}
Given a \tool{} $\Hcal(M, E)$ and a set $m_{start} \subseteq M$, the criticality of a material $m_i$ is defined as, 
\begin{align}
\nonumber
c(m_i, m_{start}, \Hcal) = |R(m_{start}, \Hcal)| - |R(m_{start}\setminus m_i, \Hcal \setminus m_i)|,
\end{align}
i.e., it is the reduction in the 
size of the reachable set when $m_i$ is removed.
\end{definition}

In our analysis of criticality, we always set $m_{start} = m_{base} \cup \mathtt{Pre}$. For each \tool{}, we compute $c(m_i, 
m_{start}, \Hcal) \forall m_i \in M$ (except for $m_{base}$ itself). Table~\ref{table:criticality} shows
the top 3 most critical materials for multiple \tool{}s. For a \tool{}, its \textbf{criticality list} is the set $M$ of 
its materials, ordered from most critical to least critical. Note that this list is partially ordered, as several materials can
have the same criticality.

\begin{table}[h]
%\vspace{-0.3in}
%\centering
\caption{Top 3 critical materials in each \tool{}. In each case, $m_{base}$ is written in bold. The criticality value
for $m_{base}$ is $R(m_{base} \cup \mathtt{Pre}, \Hcal)$.}
\begin{center}
% \begin{tabular}{|r|c|} \hline
\begin{tabular}{|R{0.8\linewidth}|l|} \hline
\toprule
Material [HS Code] & Criticality \\ 
\midrule
\textbf{Elemental Boron [284520]} & \textbf{73} \\
Hydrogen (byproduct) [280410] & 12 \\
Oxygen [280440] & 9 \\
Fumed silica [281122] & 8 \\
% Soda Ash (Sodium Carbonate) [283620] & 8 \\
% Silicon tetrachloride [281230] & 7 \\
\midrule
\textbf{Cobalt ore and concentrates [260500]} & \textbf{61} \\
Sodium Chloride (in brine solution) [250100] & 8 \\
Sodium Hydroxide (caustic soda) [281512] & 6 \\
Water (distilled or deionized) [285100] & 5 \\
% Sulfuric acid [280700] & 5 \\
% Hydrogen (elemental, compressed or liquefied) [280410] & 5 \\
\midrule
\textbf{Gallium in ores and concentrates [260600]} & \textbf{53} \\
Lithium bromide [282759] & 7 \\
Lithium chloride [282720] & 6 \\
Carbon monoxide [281121] & 5 \\
% Natural gas [271121] & 5 \\
% Methyllithium [284530] & 4 \\
% \midrule
% \textbf{Quartz (natural silica, silicon ore) [250610]} & \textbf{18} \\
% Silicon (not less than 99\% purity) [280469] & 6 \\
% Sodium silicates [283911] & 5 \\
% Silica sand [250510] & 5 \\
% Silicon dioxide (synthetic silica gel) [281122] & 4 \\
% Ferro-silico-chromium [720250] & 3 \\
\midrule
\textbf{Germanium metal, unwrought [811292]} & \textbf{40} \\
Coke (coal-derived, industrial grade) [270400] & 6 \\
Calcium carbide, industrial grade [284910] & 6 \\
Sawdust or wood shavings [440149] & 5 \\
% Carbon monoxide [282410] & 4 \\
% Lime (Calcium oxide, quicklime) [252210] & 4 \\
\bottomrule
%\hline\end{tabular}
\end{tabular}
\end{center}
\label{table:criticality}
\end{table}

\subsection{Disruption analysis with UN Comtrade}
Using the Boron \tool{} and Comtrade data~\cite{UNComtrade2024} for 2024, we
assess how disruptions in countries dominating the export of key raw
or intermediary materials (HS 6-digit level) cascade through
downstream technologies. 
UN Comtrade~\cite{UNComtrade2024} records bilateral trade by HS code for over 200 reporters; because MEKH vertices are HS-6 codes, the join is exact.

For each HS 6-digit code in the MEKH, we identify the top supplier using UN Comtrade trade values and apply a 60\% dominance threshold, i.e. the leading exporter must account for at least 60\% of
global exports of that HS code. We then simulate a disruption scenario
in which the dominant supplier halts all exports of that material, and
evaluate the downstream effects on dependent HS codes within the Boron \tool{}. Specifically, any HS codes that rely exclusively on the disrupted
precursor are considered fully removed from the network.

The analysis 
reveals that while disruptions by Mexico (HS 252922),
Norway (271121), Turkey (284011), the United States (284020), and
China (810411) do not propagate further downstream, disruptions by
Poland (HS 280200) and Turkey (HS 284019) cause significant cascading
losses. Specifically, the loss of HS 280200 eliminates five dependent
HS codes, while the disruption of HS 284019 removes
four HS codes. Additionally, the US, which controls 100\% of HS
284520, could disrupt both 284520 and its dependent HS 285000, which
lacks substitute inputs.

Overall, the results indicate that although most disruptions are
localized, a few highly concentrated supply nodes, such as Poland’s HS
280200 and Turkey’s HS 284019, pose significant downstream risks within
the boron \tool{}.

\subsection{Use in practice}
The pipeline's outputs are used in three ways at the institute. (1)~HS-code groupings derived from MEKHs define the raw and refined material baskets of a separate US--China microelectronics vulnerability analysis (in submission). (2)~The pipeline's material-classification agent, run in a lightweight export mode without full process discovery, maps each of the 73 USGS mineral commodities to its HS-6 code bucket (426 codes in total); joined to UN Comtrade US-import values for 2017--2025, these buckets drive the trade-disruption view of an interactive dependency-network explorer used by researchers at the institute, a companion to a separate top-down technology-decomposition system (in submission). That view computes, per commodity, the import value lost when the top $k$ exporters are removed and a supplier lock-in count, following the multiplex trade-network metrics of the vulnerability analysis in (1). (3)~Full MEKHs are used directly, as in this section, for the mineral-specific cascade questions that neither the commodity buckets nor the top-down decomposition can answer. Uses (1) and (2) depend only on the HS typing; use (3) also depends on the auditable references.
